\section{Introduction}
Post-training is now the workhorse for turning pretrained LLMs into usable systems: instruction tuning, preference optimization, domain adaptation, and continual updates~\citep{cite}. This phase relies on large, rapidly assembled mixtures of human-written, logged, and synthetic data. At this scale, some noise is inevitable, and often hard to remove without also discarding valuable samples.

The consequences of such noise are hard to evaluate once absorbed into a model.
\textbf{Externally}, thorough auditing is uncommon: data checks are usually coarse, and when failures appear, the responsible training fragments are difficult to isolate.
\textbf{Internally}, the problem is structural: LLM knowledge is distributed across parameters, forming an interdependent web rather than a set of independent facts. Corrupting one piece of knowledge may not stay local; it can subtly reshape related beliefs and behaviors.

This motivates two questions that are especially relevant under noisy post-training.
First, \textbf{knowledge confidence}: confidence is often treated as a proxy for mastery, yet noise may distort the confidence--correctness link, yielding systematic overconfidence rather than occasional error.
Second, the \textbf{knowledge ripple effect}: if a wrong fact is introduced, how far does its influence travel across related knowledge? Understanding this ``blast radius'' matters for data curation, post-training safety, and knowledge editing.

We investigate how \emph{noisy knowledge}\footnote{Here, \emph{noisy knowledge} means \emph{structured factual statements that contradict a trusted ground-truth knowledge base}. We instantiate noise as false knowledge triples $(s,r,o)$ that conflict with validated facts, and inject them during post-training.} affects both confidence and neighboring knowledge.
Our approach is controlled: we inject verified contradictory triples during post-training, then probe accuracy and token-level confidence on (i) the injected facts and (ii) their 1-hop/2-hop neighbors in a validated knowledge graph, allowing systematic measurement of non-local effects.

Through systematic analysis, we observe that even limited structured noise can make the model \emph{more confident when wrong} on the corrupted facts, and this miscalibration spreads to related facts in a topology-dependent manner. These results suggest that post-training evaluation should move beyond aggregate task scores and confidence-only heuristics toward structure-aware diagnostics.

\paragraph{Contributions.} We contribute: (1) evidence that noisy post-training can systematically bias knowledge confidence~\yuji{"noisy"->"small/minor" since not necessarily noisy updating can cause it.}, (2) a characterization of ripple effects across connected knowledge\yuji{Characterization sounds vague}, and (3) a general experimental framework to study noisy knowledge injection with controlled, graph-grounded probes.\yuji{We also have a dataset.}